\documentclass[11pt]{article}
\usepackage{mathunal-base}
\mathunalfuente{serif}
\providecommand{\nota}[1]{\par\smallskip\noindent{\small\itshape\color{unalblue}#1}\par\smallskip}

\renewcommand{\examtitulo}{Segundo Parcial Teórico de Métodos Numéricos (3006907), TEMA A}
\renewcommand{\examfecha}{Lunes 22 de julio de 2019}

\begin{document}

\examheader

\vspace{4pt}
\noindent\textit{Duración: 1 hora y 50 minutos.}

\vspace{6pt}
\noindent Nombre \dotfill

\noindent N$^o$ Carné \dotfill\ Grupo \dotfill

\vspace{8pt}
\noindent\textbf{La interpretación del examen hace parte de la evaluación, por tanto no se admiten preguntas comenzada la prueba.}

\vspace{10pt}
\noindent\textbf{1. (20\%)} Consideremos dos funciones infinitamente diferenciables, $f$ definida en $[-3,3]$ y $g$ definida en $[-1,1]$. Completar:
\begin{enumerate}
\item[a.] Si aplicamos la regla de trapecio compuesta con 9 subintervalos para aproximar el valor de $\displaystyle\int_{-3}^3 f(x)\,dx$, el error que se obtiene es $\mathcal O\!\left(\underline{\hspace{2cm}}\right)$.
\item[b.] El grado de precisión o exactitud al aplicar cuadratura Gaussiana con 6 nodos es \underline{\hspace{2cm}}.
\item[c.] Si $f$ es una función impar (esto es, $f(-x)=-f(x)$) y empleamos la regla de Simpson compuesta con 4 subintervalos entonces
\[
\int_{-3}^3 f(x)\,dx\approx\underline{\hspace{4cm}}
\]
\item[d.] Si denotamos por $Q$ el polinomio de grado menor igual a 5 que mejor aproxima a la función $g$ en el intervalo $[-1,1]$ y sabemos que $\displaystyle\max_{-1\leq x\leq1}|g^{(4)}(x)|\leq8.2$, $\displaystyle\max_{-1\leq x\leq1}|g^{(5)}(x)|\leq3.4$ y $\displaystyle\max_{-1\leq x\leq1}|g^{(6)}(x)|\leq1.75$, entonces la cota para el error que se comete al aproximar la función $g$ por medio del polinomio $Q$ en $[-1,1]$ es \underline{\hspace{4cm}}.
\end{enumerate}

\nota{En adelante: trabajar con redondeo a dos dígitos decimales. Recuerde: respuestas sin procedimiento no se califican.}

\vspace{5mm}
\noindent\textbf{2.} Mediante la teoría de los mínimos cuadrados, queremos hallar la curva de ajuste de la forma $y=Ce^{Dx}$ para la nube de puntos
\[
\{(x_i,y_i)\}_{i=1}^4=\{(-2,3.1),(1,4.5),(3,5.78),(6,11.3)\}.
\]
\begin{enumerate}
\item[a.] (7\%) Indique el cambio de variables apropiado para obtener la ecuación linealizada de la forma $Y=AX+B$.
\item[b.] (8\%) Plantee el sistema de ecuaciones normales a resolver para determinar $A$ y $B$.
\end{enumerate}

\vspace{5mm}
\noindent\textbf{3.} Consideremos la función $g(x)=\dfrac{100\tan^{-1}(5x)}{8x^2+3}$ y el intervalo $I:=[-4,2]$.
\begin{enumerate}
\item[a.] (12\%) Halle el polinomio que interpola a $g$ en $I$ en los nodos $x_0=-4$, $x_1=-3$, $x_2=-1$, $x_3=2$.
\item[b.] (6\%) ¿Cuál es el error relativo cometido al aproximar $g(1.3)$ por medio del polinomio interpolante hallado en a.?
\item[c.] (4\%) Halle el coeficiente polinómico de Lagrange $L_2(x)$.
\item[d.] (8\%) Halle los nodos necesarios para construir el polinomio de grado menor igual a 3 que mejor aproxima a la función $g$ en el intervalo $I$.
\end{enumerate}

\vspace{5mm}
\noindent\textbf{4.} Considere la siguiente fórmula de cuadratura
\[
\int_{-1}^1 f(x)\,dx\approx Af\!\left(-\frac12\right)+Bf(0)+Cf\!\left(\frac12\right).
\]
\begin{enumerate}
\item[a.] (12\%) Determine los valores de $A$, $B$ y $C$ que hagan que la fórmula de cuadratura sea exacta para todos los polinomios de grado tan alto como sea posible. ¿Cuál es el grado máximo de exactitud?
\item[b.] (10\%) Emplear esta fórmula de cuadratura para aproximar la integral $\displaystyle\int_1^\infty \cos\!\left(\frac1x\right)\sqrt{1-\frac1x}\,dx$.
\end{enumerate}

\vspace{5mm}
\noindent\textbf{5.} Consideremos la siguiente nube de puntos
\[
\begin{array}{c|c|c|c|c|c}
x_k & -3 & 0 & 1 & 3 & 7\\ \hline
y_k & -4 & -3 & 0 & 2 & 10
\end{array}
\]
Queremos hallar el spline cúbico que pasa por los puntos y para ello debemos resolver el siguiente sistema de ecuaciones (de acuerdo con lo visto en clase y el texto guía)
\[
3c_0+8c_1+c_2=8, \qquad c_1+6c_2+2c_3=-6, \qquad 2c_2+12c_3+4c_4=3.
\]

\nota{Recuerde que $b_j=\dfrac1{h_j}(a_{j+1}-a_j)-\dfrac{h_j}3(2c_j+c_{j+1})$ y $d_j=\dfrac{c_{j+1}-c_j}{3h_j}$.}

\begin{enumerate}
\item[a.] (4\%) ¿Cuál es el sistema de ecuaciones (cuadrado) que se debe resolver si queremos hallar el spline cúbico con terminación parabólica?
\item[b.] (9\%) ¿Cuál es el sistema de ecuaciones (cuadrado) que se debe resolver si queremos hallar el spline cúbico extrapolado?
\end{enumerate}

\vspace{5mm}
\noindent$\bigstar$ (10\% -- Extra) Si $S$ es el spline cúbico natural para la nube de puntos dada, halle $S(x)$ para $x\in[0,1]$ teniendo en cuenta que $S''(1)=-2.71$.

\vfill
\rule{\linewidth}{0.4pt}\\[1pt]
{\scriptsize\textit{Transcripción digital --- MathUNAL}\hfill\textcolor{unalblue}{\textbf{1}}}

\end{document}
